\documentclass[7.5pt]{article}

% Imports
\usepackage[utf8]{inputenc}

\usepackage{enumitem}

\usepackage{fontawesome}

\usepackage{hyperref}

\usepackage{helvet}
\renewcommand*\familydefault{\sfdefault} %% Only if the base font of the document is to be sans serif
\usepackage[T1]{fontenc}


\usepackage{parskip}[skip=0.5] % adds space between paragraphs, and also aligns items left
\usepackage{tikz} % for shapes and fancy graphics

\usepackage{calc} % for subtracting lengths/widths

\usepackage{xparse}

% TODO: reconsider margins
\usepackage[a4paper,margin=0.45in]{geometry} %a4paper sets margins, can add showframe to display guidelines

\pagenumbering{gobble}

% Define some TiKz colors
\definecolor{spacecadet}{HTML}{1D2951}
\definecolor{slategray}{HTML}{708090}
\definecolor{darkslategray}{HTML}{2f4f4f }
\definecolor{airforceblue}{rgb}{0.36, 0.54, 0.66}

% Env. variables
\def \primarycolor   {spacecadet}
\def \secondarycolor {darkslategray}
\def \tertiarycolor  {airforceblue} % may use a color just a touch lighter

\def \titlerectmin{0.15cm}
\def \titlerectmax{0.35cm}
\def \leftcolwidth{3.10cm}

\def\code#1{\texttt{#1}}

\newcommand{\CPP}
{\code{C\nolinebreak[4]\hspace{-.05em}\raisebox{.22ex}{\footnotesize\bf ++}}}



% Section commands
\newcommand{\htop}[1]{\textbf{#1}}
\newcommand{\hsub}[1]{#1}
\newcommand{\hextra}[1]{#1}

% Item commands
\newcommand{\cvheader}[1]{
    \tikz{\fill [\primarycolor] (0,\titlerectmin) rectangle (\leftcolwidth,\titlerectmax);} \ 
    \textcolor{\primarycolor}{ \large{\textbf{#1}} } 
}

\newcommand{\cvsubtitle}[1]{
    \hspace{\leftcolwidth} \ \textcolor{\secondarycolor}{ \large{{#1}} }
}

% note-- this will override the cventry command in the moderncv package
\newcommand{\cventry}[4]{
    \begin{minipage}[t]{\leftcolwidth} 
        \small{
            \begin{flushright}
                #2\\
                \textcolor{\tertiarycolor}{\textit{#4}}
            \end{flushright}
        }
    \end{minipage} \ \
    \begin{minipage}[t]{\linewidth-\leftcolwidth} \normalsize{\textbf{#1} \\ #3} \end{minipage}
} % args - title, dates, desc., website

\newcommand{\cventrysingle}[3]{
    \begin{minipage}[t]{\leftcolwidth} \small{#2 \begin{flushleft}\textcolor{\tertiarycolor}{\textit{#3}} \end{flushleft} } 
    \end{minipage} \ \
    \begin{minipage}[t]{\linewidth-\leftcolwidth} \normalsize{{#1}} \end{minipage}
}

% Misc. Commands


\setlist[itemize]{leftmargin=5.5mm}

\begin{document}

%\parskip=0.75em

{ \begin{center}\LARGE \textbf{Noah E. Wolfe} \end{center} }
%%% Summary (tl;dr)
% currently working on this
%\cvheader{Biosketch}
%Noah \textit{(he/him/his)} is a highly motivated undergraduate student and researcher at North Carolina State University. His research interests lie primarily in theoretical astrophysics and cosmology, both in numerical and "pen-and-paper" work. Ultimately, his chief interest is using astrophysics to probe the frontiers of fundamental physics, and so he has participated in research projects using a variety of tools, including numerical simulations, "pen-and-paper" theory, and novel statistical methods, to better connecting theory to observation.

%%% Education
\cvheader{Education}

\cventry{Physics B.S. (ongoing), Mathematics B.S. (ongoing)}{Aug 2018 - May 2022}{North Carolina State University}{}

%%% Research
\cvheader{Research}

\cvsubtitle{Publications}

\cventrysingle{\textbf{Wolfe, N. E.}, Fr\"ohlich, C., et al. (2021). PUSHing core-collapse supernovae to explosions in spherical symmetry VI: Gravitational Wave Eigenfrequencies. Manuscript in preparation.}{}{}

\cventrysingle{Curtis, S., \textbf{Wolfe, N.}, Fröhlich, C., Miller, J. M., Wollaeger, R., & Ebinger, K. (2020). Core-Collapse Supernovae: From Neutrino-Driven 1D Explosions to Light Curves and Spectra. \href{https://arxiv.org/abs/2008.05498}{\color{\tertiarycolor}{arXiv preprint  arXiv:2008.05498}}.}{}{}


\iffalse
\cvsubtitle{Presentations}

\cventrysingle{\textbf{Wolfe, N. E.}, Curtis, S., Ghosh, S., Fr\"ohlich, C. “Characterizing Gravitational Wave Signals from Core-Collapse Supernovae,” McCormick Undergraduate Research Symposium, Raleigh, NC, May 2020 (Talk) \href{https://drive.google.com/file/d/1l_uCcEURbu1l5_SZtgUtiaaeLIREYsWq/view?usp=sharing}{\color{\tertiarycolor}{Slides}}}{}{}

\cventrysingle{\textbf{Wolfe, N. E.}, Curtis, S., Ghosh, S., Ebinger, K., Fr\"ohlich, C. “Characterizing Gravitational Wave Signals from Core-Collapse Supernovae,” American Physical Society Division of Nuclear Physics, Crystal City, VA, October 2019 (Poster)}{}{}

\cventrysingle{\textbf{Wolfe, N. E.}, Curtis, S., Ghosh, S., Fr\"ohlich, C. “Characterizing Gravitational Wave Signals from Core-Collapse Supernovae,” Fifty-One Ergs, Raleigh, NC, May 2019 (Poster)}{}{}

\cventrysingle{\textbf{Wolfe, N. E.}, Kashani, S., Stuard, S., Gadisa, A., Ade, H. "Extraction of Exciton Binding Energy from the Stark Effect," McCormick Undergraduate Research Symposium, Raleigh, NC, April 2019 (Poster)}{}{}
\fi 

% Projects
\cvsubtitle{Projects}

\cventry{Characterizing Gravitational Wave Signals from Core-Collapse Supernovae}{Oct 2018 - Present}
{\vspace*{-2.5mm} \begin{itemize}
        \item{Modified and extended 1D core-collapse models written in \code{FORTRAN}.}
        \item{Modified and extended \code{Python} code to calculate gravitational wave eigenfrequencies given hydrodynamic data from supernovae models.}%%clown, glad u found this%%
        \item{Applied novel methods for uncertainty quantification \& sensitivity analysis, including hierarchical regression and active subspace analysis, using \code{Python} and \code{Julia}.}
        \item{"Poster-Blitz Advertisement" \& poster presentation at Fifty-One Ergs (FOE) in May 2019}
        \item{Created a system in \code{bash} to automatically run and manage a suite of models.}
        \item{Writing the final manuscript, with an anticipated publication date in Spring 2020.}
    \end{itemize}
}{with Associate Professor Carla Fr\"ohlich, NC State}

\cventry{Investigating Neutral Hydrogen in the Cold Accretion Disk around Sagittarius A*}{June 2019 - October 2020}{\vspace*{-3.5mm} \begin{itemize}
    \item Studied the statistical mechanical description and observational consequences of the $\mathrm{H}30\alpha$ transition of hydrogen from $n = 31$ to $n = 30$, in the environment around a black hole, in the presence of significant \textit{non-blackbody} background radiation.
    %\item Applied statistical mechanics, in a "pen-and-paper" manner, to understand how transition probabilities are modified by the non-blackbody background.
    %\item Independently began this collaboration with Dr. Murchikova, during a conversation following her presentation at NC State.
\end{itemize}}{with Dr. Elena Murchikova, Institute for Advanced Study}

%%\cventry{Gravitational Wave Waveforms from 1-D Models of Core-Collapse Supernovae}{May 2019 - Present}
%%{Using the characterization of the gravitational wave signals from my work with Professor Fr\"ohlich, I am working to produce gravitational wave signal "templates" that will make it easier to detect these signals from core-collapse supernovae. This project grew out of discussions at Fifty-One Ergs in May 2019.}{with Postdoctoral Fellow Sarah Gossan, CITA}

\cventry{Applying Sensitivity Analysis Methods to Core-Collapse Supernovae Lightcurves}{June 2020 - September 2020}{\vspace*{-3.5mm} \begin{itemize}
    \item Applied a novel method for sensitivity analysis, known as active subspace analysis to analyze the lightcurves resulting from 1D core-collapse supernovae models; published in Curtis et al. (2020).
    %\item Identified the produced $^{56}\textrm{Ni}$ mass as the primary determinant of lightcurve properties.
    %\item Contributed a subsection and appendix as second author to Curtis et al. (2020), in the Astrophysical Journal.
\end{itemize}}{arXiv:2008.05498}

\cventry{Machine Learning Parameter Inference with Core-Collapse Supernovae Models}{June 2020 - Present}{\vspace*{-3.5mm} \begin{itemize}
    \item Small, multi-institution collaboration using machine learning and modern statistical methods to constrain key physical parameters like the nuclear equation of state.
    \item Assist in finding astronomical data that can be compared to 1D core-collapse models.
\end{itemize}}{with Prof. Frohlich, Dr. Jonah Miller (LANL), Dr. Ricardo Vilalta (Univ. of Houston)}

%\cventry{Applying Machine Learning to Core-Collapse Supernovae Models}{}{}{}

\iffalse
\cventry{Optical Simulation of Organic Photovoltaic Devices}{Aug 2018 - July 2019}{\vspace*{-3.5mm} \begin{itemize}
    \item Created optical simulations in Python to model light transmission through organic photovoltaic devices.
    \item Attempted to calculate the change in exciton binding energy when exposed to an external electric field (via the Stark Effect).
\end{itemize}
}{with Research Professor Abay Dinku, NC State}
\fi

\iffalse
% add link to something that shows Scivir got the sustainability grant
\cventry{NC State Office of Undergraduate Research Award}{April 2020}{Awarded \$2,000 to perform research with Dr. Elena Murchikova at the Institute for Advanced Study during the summer of 2020. This travel was ultimately cancelled due to the COVID-19 pandemic.}{}

\cventry{NC State Sustainability Fund Grant}{April 2020}{Awarded \$5,820 to deploy a network of outdoor particulate matter sensors across NC State's campus, from July 2020 to June 2021. This funding is still active despite the COVID-19 pandemic.}{}
\fi

%\vspace{250px}

%%% Service
\cvheader{Service \& Leadership}

\cventry{\underline{President} \& \underline{Co-Founder}, Scivir \textit{(see-ver)}}{January 2020 - Present}{\vspace*{-3.5mm} \begin{itemize}
    \item Scivir is a 501(c)(3) nonprofit organization, with a mission to create new technologies to accelerate environmental justice.
    \item Ideated Scivir's key technical goal, to develop low-cost mesh networks of particulate matter monitors to address critical data gaps in air quality data in rural regions.
    \item Lead the organizational and technical vision of Scivir.
    \item Lead firmware development, including PCB design and \CPP{} development.
    \item Establish partnerships with communities, businesses, universities, and funding-agencies.
\end{itemize}}{scivir.org}

\cventry{\underline{President}, Astronomy Club at NC State}{May 2019 - Present}{Lead the club towards a focus on science outreach; led multiple outreach events in Raleigh, and additional remote outreach events during the fall 2020 semester.}{}

%\cventry{\underline{Co-Chair}, Park Scholars Class of 2022 Legacy Committee}{Sep 2018 - Present}{\vspace*{-3.5mm} \begin{itemize}
    %\item Lead the definition and implemention of a class legacy for the Park Scholars Class of 2022.
    %\item Lead our partnership with the Interfaith Food Shuttle in Raleigh, that harnesses the wide-ranging interests and expertise of our class towards addressing problems of food insecurity, education, and sustainability.
%\end{itemize}}{}

% Check these dates-- also steal title and desc. from prev. cv
%\cventry{Low-Cost Air Quality Monitor Mesh Networks}{June 2015 - Present}{Designing and building low cost air quality monitor mesh networks, aimed at collecting data in southeastern North Carolina; a region plagued by severe environmental justice issues caused by air pollution, where long-term air quality data does not currently exist.}{go.ncsu.edu/aq-cei \\ code @ git.io/fjQy5}

%\cventry{Small-Group Leader, Triangle Youth Leadership Conference}{January 2020}{\vspace*{-3.5mm} \begin{itemize}
    %\item Volunteer small-group leader at the Triangle Youth Leadership Conference, which brings high school students from across North Carolina to an intensive, two-day conference wherein they learn effective leadership in the context of solving community issues. 
    %\item Encouraged bonding, teamwork, and the development of critical leadership skills among a small student group, and guided these students in the development and presentation of a service prototype project.
%\end{itemize}}
%{triangleleadership.com}

%\cventry{Alternative Service Break - Alaska}{Spring Break (March) 2020}{\vspace*{-3.5mm} \begin{itemize}
    %\item Traveled to the Tlingit-majority community of Hoonah, Alaska, to understand the roots of social injustice in Native American communities. 
    %\item Volunteered at the local school and Boys \& Girls Club in Hoonah.
    %\item Learned about the intersection of social justice and educational environments relevant to being an effective mentor for students in a university setting.
%\end{itemize}}{}

%\cventry{Quantum Information Club}{August 2020 - Present}{Member of the Quantum Information Club at NC State, to learn more about quantum computing, and applications both to the wider world and my own research interests.}{}

%\cvheader{Outreach Events}

%\cventry{Athens Drive Magnet High School}{}{}{}

\cvheader{Awards \& Honors}

\cventry{Goldwater Scholarship}{April 2020}{The Goldwater Scholarship is one of the oldest and most prestigious national STEM scholarships, and identifies students with exceptional promise of becoming next generation research leaders.}{lnnk.in/@goldwater-interview}

\cventry{Park Scholarship}{Aug 2018 - Present}{NC State's highly selective, full merit scholarship awarded based on outstanding accomplishments and potential in scholarship, leadership, service, and character.}{park.ncsu.edu}

%\cventry{2nd Place - McCormick Undegraduate Physics Research Symposium}{May 2020}{Won 2nd place for my three-minute, three-slide presentation of my work Characterizing Gravitational Waves from Core-Collapse Supernovae at the NC State McCormick Undergraduate Physics Research Symposium.}{}

%\cventry{Phi Beta Kappa}{May 2020}{Invited to join Phi Beta Kappa, the oldest and best known society for recognition of academic excellence and scholarly achievement in the United States. Election to membership in Phi Beta Kappa is one of the highest honors that students can earn at NC State University.}{pbk.org}

%\cvheader{Skills}

\iffalse
\begin{description}
    \item{ 
    \textbf{General Technical Skills:} general Unix/Linux proficiency (Ubuntu, CentOS, Raspbian), shell scripting (Advanced), \LaTeX (Intermediate)
    }

	\item{
	\textbf{Scientific Computing:}
	Python (Advanced, incl. Scipy and NumPy),  Mathematica (Intermediate), Julia (Intermediate) MatLab (Intermediate), MPICH (Fundamental), Machine Learning with Python Keras library (Fundamental)
	}
	
 	\item{ 
 	\textbf{Firmware Development:} C and \CPP \ (Intermediate), development for Arduino Uno and ESP32 microcontrollers (Intermediate), use of digital sensors (Intermediate), soldering (Intermediate), Autodesk EAGLE (Fundamental)
 	}
	
	\item{
	\textbf{Web Development}: HTML and template engines incl. ejs (Advanced), JavaScript (Advanced), node.js (Advanced)
	}
	
	\item{
	\textbf{Data Systems/Formats:}
	HDF5 (Advanced), PostgreSQL (Intermediate), MongoDB (Fundamental)
	}
	
	\item{
	\textbf{Other Technical Experience:} Soldering/electronics (Intermediate), Arduino Uno, Raspberry Pi, ESP32
	}
	
	\item{
	\textbf{Languages:} English (native), Gujarati (Fundamental), Spanish (Fundamental)
	}
	
\end{description}
\fi

\iffalse
%%% Biosketch -- maybe put this on it's own page
Noah Wolfe is an undergraduate student and Park Scholar pursuing degrees in both Physics and Mathematics at North Carolina State University (NC State). His work thus far has combined his talent for programming with his love of astrophysics, as he has studied core-collapse supernovae, neutron stars, and gravitational waves with state-of-the-art simulations. However, his astrophysical research interests are not limited to the theoretical; he has recently become involved in a project to study the environment surrounding Sagittarius A* with ALMA observations. Noah also strongly believes in serving his communities. As President of the Astronomy Club, he has lead scientific outreach events both on-campus and in the greater Raleigh community. He has also lead the development of low-cost air quality sensor networks, in order to collect air quality data in rural, socioeconomically disadvantaged regions of North Carolina where this data does not currently exist. In his free time, Noah enjoys playing the drum set, baking vegan desserts, hiking, and reading sci-fi novels. 
\fi

\end{document}
